% CSC 5323 "Advances" paper: literature review and pilot of logq.
% Template: ICSE 2027 Research Track requires the IEEE conference proceedings
% template, \documentclass[10pt,conference]{IEEEtran} without compsoc/compsocconf;
% 10 pages of main text plus 2 pages of references; double-anonymous review.
% Source: https://conf.researchr.org/track/icse-2027/icse-2027-research-track (read 2026-11-15).
% Draft prose by Claude Opus 5.5 (claude-opus-5-5); see the AI Assistance Disclosure
% section and ../search-log.md. Pilot numbers come from the TeamMatch write-up
% (teammatch 3616eb4:docs/advances/logq-pilot.md). Every remaining \todo{} is for the
% author; judgments marked "user to confirm" must be the author's own.
\documentclass[10pt,conference]{IEEEtran}
\usepackage{cite}
\usepackage{url}
\usepackage{xcolor}
\usepackage{array}
\newcommand{\todo}[1]{\textcolor{red}{[TODO: #1]}}
\newcommand{\cmd}[1]{\texttt{#1}}

\begin{document}

\title{Filtered Runtime Context for Coding Agents:\\A Literature Review and a Pilot of logq}
% Double-anonymous for a real ICSE submission; for the course, put your name here.
\author{\IEEEauthorblockN{\todo{author}}\IEEEauthorblockA{Louisiana Tech University}}
\maketitle

\begin{abstract}
AI coding agents debug software by running commands and reading what comes back. When an application runs as several Docker containers, the evidence is in container logs and request traces: output that is long, interleaved across services, sometimes secret, and partly written by whoever sends the application a request. We review six studies that bear directly on giving an agent such runtime context: interfaces designed for agents, agents that fetch telemetry, how models use long inputs, and the risks of logs, injected instructions and tool protocols. We then describe logq, a command-line tool and agent skill that gives a coding agent filtered, redacted, time-ordered logs across a Docker Compose project, a live filtered log stream, and request traces through Jaeger's Model Context Protocol (MCP) server, and we pilot it on TeamMatch, a Django course project. In a pilot experiment, Claude Opus 5.5 in Claude Code debugged three seeded faults once with the Docker CLI alone and once with logq and Jaeger. It found every root cause in both conditions; with logq it used 27 rather than 30 turns, about 20\% less API-priced cost and a third fewer characters of logs and traces, almost all of the difference coming from one task. With one run per cell, this is a trend, not proof.
\end{abstract}

\section{Introduction}

LLM agents now edit repositories, run tests and resolve real GitHub issues through tools, and the design of the interface between agent and computer measurably changes how well they do~\cite{yang2024sweagent}. Debugging a running application is harder to hand them. When the software is a set of containers (an API, a worker, a database, object storage, a web front end), the evidence of a failure is spread across each container's output: a stack trace in one service, a retried call in another, an exit code from a container that crashed without printing anything. A coding agent typically reads this with \cmd{docker logs}, one container at a time. That floods its context with unrelated lines, cannot connect one request's lines across services, and misses crashes that leave no log line. It also hands the model text that it uses poorly when the relevant part sits in the middle~\cite{liu2024lostmiddle}, that contains personal data~\cite{aghili2025logprivacy}, and that can act as instructions to it~\cite{greshake2023injection}.

The advance studied here is to give a coding agent \emph{filtered, redacted, live runtime context} instead of raw logs. Our tool, logq, reads container logs through the Docker daemon on demand, filters them by service, level, trace ID, route or status, merges them in time order, removes secrets, and reports containers that exited or are unhealthy. A \emph{follow} mode streams only matching lines and lifecycle events while a user reproduces a problem, and the agent reads each request's flow across services from a local Jaeger through its MCP server~\cite{hou2026mcp}. The approach is immature: the Jaeger release whose documentation first describes its MCP server appeared two days before our design was written, and the agent we found that fetches logs as tools targets incident response in a production cloud~\cite{roy2024rcaagents}.

This paper makes three contributions: a focused review of the six studies that logq's design and evaluation rest on (Section~\ref{sec:related}); the design of logq and its agent skill (Section~\ref{sec:approach}); and a pilot in TeamMatch, a course team project, with a debugging experiment (Section~\ref{sec:eval}) and its threats to validity (Section~\ref{sec:threats}).

\section{Related Work}\label{sec:related}

The sources come from an AI-assisted search run with Claude Opus 5.5 (\cmd{claude-opus-5-5}) and a web search tool. It admitted peer-reviewed papers and arXiv preprints from 2021 to 2026, mainly from ICSE, FSE, ISSTA, NeurIPS and ICLR, and widened to other venues where the canonical paper on a needed topic appeared; every candidate was checked against its DOI, arXiv record or proceedings page. That search produced 25 verified candidates across five topics: agents for debugging and repair, agent interfaces, runtime context for code models, LLMs over logs and telemetry, and tool protocols and their risks. This paper keeps the six whose findings a logq design decision or our evaluation method rests on directly. Where two candidates supported the same point, we kept the one closer to logq's setting. Every query, criterion, refinement and the full list of 25 are in the accompanying \cmd{search-log.md} and \cmd{source-matrix.md}. Table~\ref{tab:sources} lists the six, how each was evaluated, and what this paper uses it for.

\emph{Interfaces designed for agents.} SWE-agent argues that language-model agents are a new category of end user with their own needs, and that they benefit from interfaces built for them as humans benefit from IDEs~\cite{yang2024sweagent}. Its agent--computer interface (ACI) gives the agent commands to edit files, navigate a repository and run tests; with it the agent reached 12.5\% pass@1 on SWE-bench and 87.7\% on HumanEvalFix, and the paper analyzes how interface design changes agent behavior. The argument transfers to runtime output: \cmd{docker logs} is built for a human with scrollback and \cmd{grep}, while an agent with a bounded context needs output that is already filtered, ordered and summarized, and that says what it did not show. SWE-agent is evaluated by holding the model and tasks fixed and varying the interface, which is how we evaluate logq.

\emph{Agents that fetch telemetry.} Roy et al.\ observe that earlier LLM approaches to root-cause analysis could not collect additional diagnostic data such as logs, metrics or databases. They evaluate a ReAct agent whose retrieval tools fetch such data, on out-of-distribution production incidents at Microsoft: it was competitive with strong baselines with much higher factual accuracy, and a case study gave the agent a team's own diagnostic services as tools~\cite{roy2024rcaagents}. logq gives a coding agent the same ability, but in a different setting: a developer debugging local containers during development, with no telemetry pipeline, rather than on-call engineers diagnosing a production cloud.

\emph{Why filter.} Liu et al.\ place the relevant item at different positions in multi-document question answering and key-value retrieval: models use information at the start or end of their input best, degrade when it sits in the middle, and degrade further as inputs grow, including models built for long contexts~\cite{liu2024lostmiddle}. A day of interleaved container logs is such an input, which is why logq filters before the agent reads.

\emph{Risks of handing logs to an agent.} Greshake et al.\ show that LLM-integrated applications blur data and instructions: instructions planted in data the application retrieves can steer its behavior and the APIs it calls, demonstrated against Bing Chat and code-completion engines~\cite{greshake2023injection}. Logs are such data, since anyone who can make a service log a string, for example through a request path or header, can put text in front of the agent. Aghili et al.\ find personally identifiable information and quasi-identifiers in 25 public log datasets, review what GDPR and CCPA require, and survey 45 practitioners, concluding that the field lacks consistent definitions of sensitivity and standard guidelines~\cite{aghili2025logprivacy}. Hou et al.\ define an MCP server's lifecycle in four phases and 16 activities, derive a taxonomy of four attacker types and 16 threat scenarios, and propose safeguards for each phase~\cite{hou2026mcp}; reading traces through Jaeger's MCP server inherits that threat surface.

\emph{Gap.} The agent that fetches logs targets production incidents, the interface work targets code and tests, and the risks of logs are studied through attacks, datasets and surveys rather than in an agent evaluation. None of the six evaluates an agent debugging a development deployment from filtered runtime output, or reports whether secrets in that output reached the model. logq and its pilot address both, at small scale.

\begin{table}[t]
\caption{The six sources this paper keeps}\label{tab:sources}
\centering\footnotesize\setlength{\tabcolsep}{4pt}
\begin{tabular}{>{\raggedright}p{0.2\columnwidth}>{\raggedright}p{0.34\columnwidth}>{\raggedright\arraybackslash}p{0.34\columnwidth}}
\hline
\textbf{Work} & \textbf{Evaluation} & \textbf{Used here for} \\
\hline
SWE-agent~\cite{yang2024sweagent} & SWE-bench, HumanEvalFix; interface varied & agent-facing output; evaluation method \\
Roy et al.~\cite{roy2024rcaagents} & Production incidents; team case study & closest agent that fetches logs as tools \\
Liu et al.~\cite{liu2024lostmiddle} & Retrieval tasks, relevant item moved & filtering before the agent reads \\
Greshake et al.~\cite{greshake2023injection} & Attacks on deployed LLM applications & log content as data, never instructions \\
Aghili et al.~\cite{aghili2025logprivacy} & 25 log datasets; laws; 45 practitioners & always-on redaction \\
Hou et al.~\cite{hou2026mcp} & MCP lifecycle and threat taxonomy & risks of Jaeger's MCP server \\
\hline
\end{tabular}
\end{table}

\section{Approach: logq}\label{sec:approach}

When an agent debugs an application in local Docker containers, it should see only the log lines that matter, from every container, in time order, with secrets removed. Concretely: a query for one trace ID returns that request's lines from every service, oldest first; every query states what it searched, so an empty result is never mistaken for success; failed containers are reported first; and no secret pattern from the test fixtures appears in any output.

logq is a Python command-line program with no dependencies beyond the standard library and no collector: the Docker daemon already stores and labels every container's logs, so logq reads them through the \cmd{docker} CLI only while a command runs, under the user's account and Docker context. Each command discovers a Compose project's containers from their labels, reads their logs in parallel, parses each line into an event, filters, then redacts, folds multi-line errors, caps and prints.

\cmd{logq ps} lists each container with its state, health, uptime and restarts, marking exited, restarting and unhealthy ones. \cmd{logq query} reads a time window (default 15 minutes) and prints the most recent matches (default 50) oldest first, filtered by service, minimum level (default \emph{warn}), trace or request ID, route, method, status class, code function or a regular expression. A footer names the containers, project, window and level, says when matches were dropped, and counts events that matched every filter but fell below the level (Fig.~\ref{fig:query}). \cmd{logq follow} applies the same filters live, adds Docker's container events (die with exit code, OOM, restart, start, unhealthy) regardless of level, and collapses repeats within 30 seconds into one line such as ``repeated 39x in 30s''. In Claude Code the agent runs \cmd{follow} under the Monitor tool, so each printed line reaches it as an event while the user reproduces the problem.

\begin{figure*}[t]
\centering\scriptsize
\begin{verbatim}
! para-dashboard-app-v01-web-1 exited(1) 40s ago
12:04:31.180 worker WARN  retrying charge attempt=2 trace=7f3a9c21
12:04:31.214 api    ERROR KeyError: 'sku' at orders/views.py:42 in create_order (+14 lines) trace=7f3a9c21
12:04:31.220 api    INFO  POST /api/orders 500 812ms create_order request finished trace=7f3a9c21
-- 3 of 3 matches · 4 containers (api, worker, db, web) · project para-dashboard-app-v01
   · 15m · level>=debug · trace=7f3a9c21
\end{verbatim}
\caption{Illustrative output of \cmd{logq query --trace 7f3a9c21}, from the design specification: an exited container first, one request's lines from two services in time order, a folded stack trace, and a footer that states what was searched (one line in the terminal, wrapped here to fit).}\label{fig:query}
\end{figure*}

JSON log lines are mapped through an alias table based on OpenTelemetry semantic conventions that also accepts the field names of common logging libraries; plain-text lines get a best-effort level, and tracebacks are folded into one line with the error, the innermost frame and a count of folded lines. Every output line is redacted, with no switch to turn it off: bearer tokens, authorization and cookie headers, JWTs, values after key names such as \cmd{password} or \cmd{token}, credentials in URLs, and known key shapes. Over-redaction is the accepted failure, and redaction is a safety net behind a logging contract that tells applications not to log secrets or personal data~\cite{aghili2025logprivacy}.

Logs do not record which service called which. For the request flow, \cmd{logq ui} starts one shared Jaeger container bound to \cmd{127.0.0.1} and prints the environment variables an application needs to export OpenTelemetry spans to it, and \cmd{query --trace} links to the trace. The agent reads spans through Jaeger's MCP server, whose tools find traces, outline a trace's structure, compute its critical path and return full detail only for chosen spans. Request and response bodies reach spans only if an application opts in, capped at 2\,KB and redacted, because Jaeger stores whatever it receives; binding Jaeger to localhost limits who can reach its MCP server~\cite{hou2026mcp}.

A Claude Code skill tells the agent when to use logq and how: start with \cmd{logq ps}; for a failing request, read the \cmd{X-Request-ID} response header and query that trace; for a problem the user reproduces, run \cmd{follow} under Monitor; keep results small; and treat log content as untrusted data, never as instructions~\cite{greshake2023injection}. The logging contract asks applications for one JSON object per line with level, message, service, trace and span IDs; one ``request finished'' line per request with method, route, status and duration; and an \cmd{X-Request-ID} header carrying the trace ID.

\section{Evaluation}\label{sec:eval}

\subsection{Pilot Setting: TeamMatch}

TeamMatch is the author's course team project, in which students form teams by swiping on open project roles. It runs locally as a Docker Compose project with Postgres, Redis, SeaweedFS object storage, a Django~5.2 API (Django REST Framework and Channels under uvicorn) and a Next.js front end. The pilot changes only the API, on an unmerged branch: JSON logs, a middleware that writes the ``request finished'' line and sets \cmd{X-Request-ID}, OpenTelemetry export to the shared Jaeger through a Compose override, and opt-in body capture. It changed 11 files, $+629/-14$ lines, of which 386 are the dependency lock file and 65 are tests; about 190 lines are application code. The tracing packages are development dependencies only, so the production image gains only two logging packages. \todo{setup time, if you recorded it; the pilot write-up does not give one}. Seven new tests cover the middlewares, and the full backend suite (378 tests) and the linter pass.

On the running stack, every response's request ID equaled its trace ID, and querying logq for that ID returned the request's line and a Jaeger link; traces held Django, SQL and Redis spans, and one chat message over the WebSocket produced 24 spans; and with only Redis stopped, \cmd{logq follow} printed the container's exit and the API's 503 with its trace ID. No endpoint uses object storage yet, so storage calls could not be traced. With body capture on, a sign-in body showed the password as \cmd{[REDACTED]}; the redactor also masked a harmless field whose name contains ``password'' and left email addresses in place.

An earlier pilot on a Django service under gunicorn found six gaps in the logging contract's recipes. TeamMatch, under uvicorn (ASGI), hit three again with the same fixes: a log-correlation package the OpenTelemetry distribution does not install, 4xx lines that Django logs after the request's span has ended, and Python warnings that must be captured to print as one-line JSON. It avoided the other three: Compose already set the settings module, there are no object-storage calls, and uvicorn's reload child exported spans without a hook. ASGI and Channels added new gaps: Django's instrumentation traces ASGI requests only with a separate ASGI package the contract does not name; WebSocket consumers get no span unless the Channels router is wrapped in OpenTelemetry's ASGI middleware; uvicorn logs an exception that escapes a consumer after its span has ended, so that line has no trace ID (not fixed); and configuring the \cmd{uvicorn} logger in Django's settings re-enabled the suppressed access log. Jaeger~2.23 also no longer serves two of its v1 search endpoints, while fetching a trace by ID still works.

\subsection{Research Questions and Design}

\begin{itemize}
\item \textbf{RQ1 (effectiveness).} With logq, does the agent identify the root cause more often than with the Docker CLI alone?
\item \textbf{RQ2 (efficiency).} Does logq change the turns, tokens, time and tool output the agent spends on the way to the root cause?
\item \textbf{RQ3 (experience).} Is the workflow effective, cumbersome, accessible, and appropriate for this project, and when would we recommend it?
\end{itemize}

As in the ACI evaluations~\cite{yang2024sweagent}, two conditions share the model and tasks and differ only in runtime-context tooling. In the \textbf{baseline}, the agent reads logs with the Docker CLI alone; running logq, reaching Jaeger's ports and using any MCP server are denied. In the \textbf{logq} condition, the skill's text is appended to the system prompt and Jaeger's MCP server is allowed. In both, the agent may run shell commands and read and search files, but may not edit files, browse the web, start subagents, or stop, restart, remove or recreate containers; a dry run with a smaller model checked these limits first. Each run is a fresh headless session of Claude Code 2.1.317 with Claude Opus 5.5 (\cmd{claude-opus-5-5}), capped at 40 turns, with a clean environment and no user settings, hooks or plugins. Each fault lives in its own copy of the source, exported without git history so \cmd{git diff} cannot reveal it. Before each run only the API container was recreated from that copy, keeping the data volumes, and a script signed in a seed account and triggered the symptom once. The prompt is identical apart from the symptom: it names the stack, the Compose project and helper scripts, quotes the symptom as a user would report it without the cause, and asks for one line naming the root cause as \emph{file:line} and one sentence proposing a fix. There are three tasks, each run once per condition (six runs), with the logq run first. All three are seeded in the API:
\begin{itemize}
\item \textbf{T1, feed error.} A changed guard in the feed serializer lets a role with no required skills divide by zero. Symptom: ``The swipe deck won't load for students any more; it shows an error instead of cards.'' Evidence: a 500 response, a \cmd{ZeroDivisionError} traceback in the API log, and an error span.
\item \textbf{T2, slow feed.} The feed's query no longer prefetches related rows, an N+1 pattern. Symptom: the swipe deck ``got noticeably slower after the last merge. Nothing errors, it is just slow.'' Evidence: in traces only; the request still returns 200, but its SQL queries rise from 6 to 95 and its latency from 48 to 171\,ms.
\item \textbf{T3, chat drops.} The Channels Redis \cmd{socket\_timeout} drops from 15\,s to 1\,s, below the channel layer's 5\,s blocking read. Symptom: team chat ``keeps disconnecting'', and messages reach teammates only after a refresh. Evidence: the socket closes about 1.2\,s after connecting, and the API log shows a Redis \cmd{TimeoutError}.
\end{itemize}
A planned object-storage fault was replaced by T2 because no endpoint uses storage.

\emph{Measures}, all from Claude Code's machine-readable transcript of each run: \emph{success}, when the final line names the file and cause of the seeded change (the proposed fix is recorded but not applied or tested); \emph{turns and tool calls}; \emph{tokens}, with uncached input, cache writes, cache reads and output reported separately; \emph{cost}, Claude Code's API-price estimate, since the runs used a subscription plan and were not billed; \emph{wall-clock time}; \emph{context volume}, the characters returned by log commands (logq or \cmd{docker logs}) and Jaeger's MCP tools, and of all tool output including source files; and \emph{secret exposure}, whether a seed account's password appears anywhere in the transcript.

\subsection{Results}

\begin{table*}[t]
\caption{Pilot results, one run per task and condition. Tokens and characters are counts; cost is Claude Code's API-price estimate (not billed). Per-run costs are rounded, so they may not sum exactly to the totals.}\label{tab:results}
\centering\footnotesize\setlength{\tabcolsep}{4pt}
\begin{tabular}{llcrrrrrrrrrr}
\hline
\textbf{Task} & \textbf{Condition} & \textbf{Root} & \textbf{Turns} & \textbf{Tool} & \textbf{Wall} & \textbf{Input} & \textbf{Cache} & \textbf{Cache} & \textbf{Output} & \textbf{Cost} & \textbf{Log+trace} & \textbf{All tool} \\
 & & \textbf{cause} & & \textbf{calls} & \textbf{(s)} & & \textbf{write} & \textbf{read} & & \textbf{(USD)} & \textbf{chars} & \textbf{output chars} \\
\hline
T1 feed error & Baseline & found & 8 & 7 & 27 & 16 & 30,274 & 165,952 & 1,986 & 0.315 & 9,391 & 18,358 \\
 & logq & found & 9 & 8 & 32 & 18 & 15,143 & 227,855 & 2,130 & 0.209 & 5,679 & 14,674 \\
T2 slow feed & Baseline & found & 13 & 12 & 37 & 24 & 24,634 & 375,926 & 3,164 & 0.336 & 7,836 & 45,272 \\
 & logq & found & 14 & 13 & 37 & 28 & 21,236 & 330,683 & 3,076 & 0.298 & 11,505 & 23,948 \\
T3 chat drops & Baseline & found & 9 & 8 & 31 & 14 & 16,653 & 176,671 & 2,196 & 0.213 & 16,981 & 23,817 \\
 & logq & found & 4 & 3 & 22 & 8 & 16,526 & 91,292 & 1,624 & 0.183 & 5,481 & 18,400 \\
\hline
Total & Baseline & 3/3 & 30 & 27 & 95 & 54 & 71,561 & 718,549 & 7,346 & 0.863 & 34,208 & 87,447 \\
 & logq & 3/3 & 27 & 24 & 91 & 54 & 52,905 & 649,830 & 6,830 & 0.690 & 22,665 & 57,022 \\
\hline
\end{tabular}
\end{table*}

\emph{RQ1.} The agent named the seeded file and cause in all six runs, 3/3 in each condition (Table~\ref{tab:results}), so success does not separate the conditions on these tasks. In T2, both runs pointed at the start of the feed query rather than at a line, because the removed lines no longer exist in that tree.

\emph{RQ2.} Over the three tasks, the logq condition used 27 turns against 30, an estimated cost of \$0.690 against \$0.863 (about 20\% lower), 22,665 against 34,208 characters of logs and traces (a third fewer), and 57,022 against 87,447 characters of tool output overall (35\% less). Wall time was 91\,s against 95\,s; at 22--37\,s per run, a difference of a few seconds is noise. Almost all of the gap comes from T3 (4 turns against 9): one filtered \cmd{logq query} surfaced the Redis \cmd{TimeoutError} and one search for \cmd{socket\_timeout} found the setting, while the baseline read the consumer and chat service, about 17,000 characters of raw \cmd{docker compose logs}, the settings, and the Channels Redis library's source inside the container. In T1 the conditions were close (9 turns against 8): the logq run read the traceback by querying the response's \cmd{X-Request-ID} with \cmd{--full}, and the baseline searched \cmd{docker compose logs} for the exception. In T2 traces did not shorten the search (14 turns against 13): Jaeger's MCP tools showed the repeated SELECT spans, while the baseline counted the queries itself in a Django shell, and both runs found the feed's query-count test.

\emph{Secret exposure.} No transcript in either condition contains a seed password, and no run read the file that holds the seed accounts.

\todo{RQ3: describe what the narrated screen recording shows: setup, the agent calling logq and the Jaeger MCP tools, and one task end to end.}

\subsection{Discussion}

\emph{Effective?} \todo{user to confirm: whether logq helped, on which tasks, and why, based on the results above}.

\emph{Cumbersome?} The integration cost is front-loaded: an application must follow the logging contract before logq can filter by trace or route, and the earlier pilot shows that the recipes can fail silently. \todo{user to confirm: how much effort TeamMatch took and whether it felt cumbersome}.

\emph{Accessible?} logq needs Python, the Docker CLI and, for traces, a local Jaeger container; it has no service to deploy. \todo{user to confirm: how easy it was for you and would be for teammates}.

\emph{Appropriate, and would we recommend it?} \todo{user to confirm: your recommendation, and when you would or would not use it}.

\section{Threats to Validity}\label{sec:threats}

\emph{Internal.} The author built logq, chose the tasks and graded success, which can bias results toward the tool; fixing the prompts and the success check before any run reduces but does not remove this. Each task ran once per condition, so a per-task difference could reverse in another run; repeating each cell three to five times is the next step. The logq run came first for every task, so order was not counterbalanced. The logq condition differs in more than its tools: its system prompt is about 7\,KB longer, because it carries the skill's instructions on how to debug. Both conditions had the pilot's JSON logs, so the experiment compares the tooling, not the logging contract. The agents could read the full source, including tests: in T2 the feed's query-count test pointed both runs at the missing prefetch. Log noise affects the baseline more than logq; recreating the API container before each run limits it for the API's own log.

\emph{Construct.} Turns, tokens and time are proxies for debugging effort. Token counts depend on how the platform caches input, so we report the token kinds separately, and cost is an estimate rather than a bill. Success means naming the root cause; the proposed fixes were not tested. Seeded faults have a known answer, which real debugging lacks, and the symptom text may reveal more or less than a real report. With source access and structured logs, these faults were easy for the model in both conditions, leaving little room for a difference in success.

\emph{External and conclusion.} The pilot covers one application, three faults in one service, one stack, one agent and model, and one machine, and says nothing about production systems~\cite{roy2024rcaagents}. Jaeger's MCP server is new, and MCP servers carry their own threats~\cite{hou2026mcp}. With three tasks and one run each, the results show a trend on these tasks and support no statistical claim; we report counts and totals only.

\emph{Review.} The search was AI-assisted, used a general web search tool rather than digital-library queries, and may have missed relevant work. Keeping six of 25 candidates was a judgment about relevance to logq's design and evaluation; the other 19, with the reason each was left out, are in the search log, and every cited record was checked against its DOI, arXiv record or proceedings page.

\section{Conclusion}

Prior work shows that agents do better with interfaces designed for them, that agents fetching logs and traces have so far been built for production incident response, and that putting logs in front of a model risks context overload, leaked personal data and injected instructions. logq brings these together for development-time debugging of containerized applications: filtered, time-ordered, redacted logs across a Compose project, a live filtered stream, and request traces through Jaeger's MCP server. In the TeamMatch pilot, the agent found all three seeded root causes with and without logq; with logq it used 27 rather than 30 turns, about 20\% lower cost and a third fewer characters of logs and traces, almost all of it on one task, which with one run per cell is a trend, not proof. \todo{user to confirm: your overall recommendation}. Future work includes repeated runs over more faults and services, and measuring robustness to instructions planted in logs~\cite{greshake2023injection}.

\section*{AI Assistance Disclosure}

This paper was prepared with Claude Code using the model Claude Opus 5.5 (model id \cmd{claude-opus-5-5}). The model ran the literature search and verified every source; the queries, inclusion criteria, refinements, results and verification steps are in \cmd{search-log.md}, with per-source notes in \cmd{source-matrix.md}, both beside this paper. At the student's request, the paper cites only the six most relevant of the 25 verified sources; the selection is recorded in \cmd{search-log.md}. The model also drafted the prose of every section: the Abstract, Introduction, Related Work, Approach (from the logq design specification), Evaluation, Threats to Validity, and Conclusion. The pilot's integration and experiment were also run in Claude Code sessions; the Evaluation's numbers and observations were transcribed by the model from the pilot's committed write-up in the TeamMatch repository (\cmd{docs/advances/logq-pilot.md}), and the six measured debugging runs themselves used \cmd{claude-opus-5-5}. It did not invent results; judgments marked for the author are the author's to make. \todo{user to confirm: the instructor's permission for AI-drafted sections and its conditions, and what you reviewed and changed}.

\bibliographystyle{IEEEtran}
\bibliography{../references}

\end{document}
